\section{Transformasi Fourier}\label{section_Fourier_transform}

Dalam~\ref{defn_Fourier_transform}, kita mendefinisikan transformasi Fourier suatu fungsi $\lfs 1\R$. Cara ekuivalen untuk menyatakan
definisi ini adalah melalui konvolusi:
  \[ (\mathfrak F f)(x) = (f \ast \varepsilon^x)(0) \]
untuk $x \in \R$, dengan $\varepsilon^x$ fungsi $t \mapsto e^{itx}$.

\begin{defn} Kita ingin memperluas transformasi yang berguna ini ke distribusi. Namun, ternyata ruang $\fml D^*(\R)$ terlampau besar
untuk mendukung definisi yang wajar baginya. Dengan memulai dari himpunan ``fungsi uji'' yang lebih besar, yakni fungsi-fungsi dalam ruang Schwartz
$\fml S(\R)$ (yang memuat $\fml D(\R)$ sebagai subruang rapat), kita memperoleh ruang fungsional linear kontinu yang lebih kecil; pada ruang ini, untungnya, kita
dapat mendefinisikan suatu \emph{transformasi Fourier} secara alami.
\end{defn}

Sebelum mendefinisikan transformasi Fourier pada distribusi, mari kita ingat kembali dua fakta penting dari analisis real mengenai
transformasi Fourier yang bekerja pada fungsi mulus terintegralkan.

\begin{prop}[Rumus Inversi Fourier] Jika $f \in \fml S(\R)$,
 \index{Fourier!rumus inversi}%
maka
   \[ f(x) = \int_\R \hat f \varepsilon^x\,dm\,. \]
\end{prop}

\begin{prop}\label{prop_FT_isomor1} Transformasi Fourier
   \[ \mathfrak F\colon \fml S(\R) \sto \fml S(\R)\colon f \mapsto \hat f \]
merupakan isomorfisme ruang Fr\'echet (yakni, pemetaan linear kontinu bijektif dengan invers kontinu), dan untuk semua
$f \in \fml S(\R)$ serta $x \in \R$, kita mempunyai $\mathfrak F^2f(x) = f(-x)$, sehingga $\mathfrak F^4f = f$.
\end{prop}

\begin{proof} Pembuktian kedua hasil sebelumnya, bersama banyak informasi lain tentang transformasi yang menarik ini
dan beragam penerapannya pada persamaan diferensial biasa maupun parsial, dapat ditemukan di berbagai sumber.
Beberapa yang segera terlintas adalah~\cite{Cerda:2010}, Bab~7; \cite{Horvath:1966}, Bab~4, Seksi~11;
\cite{McDonaldW:1999}, Seksi~11.3; \cite{Rudin:1991}, Bab~7; \cite{Treves:1967}, Bab~25; dan \cite{Yosida:1965},
Bab~VI.   \ns
\end{proof}

\begin{defn} Fungsi inklusi $\iota\colon \fml D(\R) \sto \fml S(\R)$ bersifat kontinu. Maka pemetaan antara ruang-ruang dualnya
   \[ u\colon \fml S^*(\R) \sto \fml D^*(\R)\colon \tau \mapsto u_\tau = \tau \circ \iota \]
merupakan isomorfisme ruang vektor dari $\fml S^*(\R)$ ke suatu ruang distribusi pada~$\R$, yang kita sebut ruang
 \index{distribusi tempered}%
\df{distribusi tempered} (sebagian penulis memilih frasa
 \index{distribusi!temperate}%
\df{distribusi temperate}). Lazimnya, fungsional linear $\tau$ dan $u_\tau$ di atas diidentifikasi. Jadi, $\fml S^*(\R)$
sendiri dipandang sebagai ruang distribusi tempered.
\end{defn}

Jika Anda kesulitan memverifikasi salah satu pernyataan dalam definisi sebelumnya, lihat Teorema 7.10\emph{ dan seterusnya} dalam~\cite{Rudin:1991}.

\begin{exam} Setiap distribusi dengan tumpuan kompak merupakan distribusi tempered.
\end{exam}

\begin{defn} Definisikan
 \index{Fourier!transformasi}%
 \index{transformasi!Fourier}%
 \index{F@$\mathfrak F(u)$ (transformasi Fourier dari~$u$)}%
 \index{<unoptop@$\hat u$ (transformasi Fourier dari~$u$)}%
\df{transformasi Fourier} $\hat u$ dari suatu distribusi tempered $u \in \fml S^*(\R)$ dengan
   \[ \hat u(\phi) = u(\hat\phi) \]
untuk semua $\phi \in \fml S(\R)$. Notasi $\mathfrak Fu$ merupakan alternatif bagi~$\hat u$.
\end{defn}

\begin{prop} Jika $f \in \lfs1\R$, maka kita dapat memperlakukan $\tilde f$ sebagai distribusi tempered dan, dalam hal ini,
   \[ \hat{\tilde f} = \tilde{\hat f}\,. \]
\end{prop}

\begin{exam} Transformasi Fourier fungsi delta Dirac diberikan oleh
   \[ \hat\delta = \tilde{\vc 1}\,. \]
\end{exam}

\begin{exam} Transformasi Fourier distribusi regular $\tilde{\vc 1}$ diberikan oleh
   \[ \hat{\tilde{\vc 1}} = \delta\,. \]
\end{exam}

Proposisi~\ref{prop_FT_isomor1} memiliki generalisasi yang memuaskan bagi distribusi tempered.

\begin{prop} Transformasi Fourier $\mathfrak F\colon \fml S^*(\R) \sto \fml S^*(\R)$ merupakan pemetaan linear kontinu
bijektif yang inversnya juga kontinu. Lebih lanjut, $\mathfrak F^4$ merupakan pemetaan identitas pada~$\fml S^*(\R)$.
\end{prop}

Kita akan kembali sejenak ke pembahasan transformasi Fourier dalam bab berikutnya.
















\endinput
